\documentclass[10pt,a4paper]{article}

\usepackage[a-1b]{pdfx}

\usepackage{fontspec}
\usepackage{geometry}
\usepackage{tabularx}
\usepackage{array}
\usepackage{etoolbox}  %% \ifblank für die Angaben, die leer sein dürfen
\usepackage{tikz}      %% das Logo

\setmainfont{FiraSans}[
  Ligatures      = TeXOff,
  Extension      = .otf,
  UprightFont    = *-Regular,
  BoldFont       = *-Bold,
  ItalicFont     = *-Italic,
  BoldItalicFont = *-BoldItalic,
]

\geometry{a4paper, left=20mm, right=20mm, top=18mm, bottom=16mm}

\setlength{\parindent}{0pt}
\pagestyle{empty}
\raggedright

\hyphenpenalty=10000
\exhyphenpenalty=10000
\tolerance=2000

\frenchspacing

\renewcommand{\arraystretch}{1.35}

\newlength{\betragsbreite}\setlength{\betragsbreite}{24mm}
\newcolumntype{B}{>{\raggedleft\arraybackslash}p{\betragsbreite}}

\newcommand{\positionstrenner}{\\\hline}
\newcommand{\summentrenner}{\\}
\newcommand{\zeilenwechsel}{\\}
\newcommand{\zusammen}[1]{\mbox{#1}}

\newcommand{\kopfangabe}[2]{\ifblank{#2}{}{#1 & #2\\}}

\newcommand{\hinweisabsatz}[1]{\par\vspace{5mm}\noindent%
  {\nonfrenchspacing\sfcode`\:=1000 \sfcode`\;=1000 \sfcode`\,=1000 #1}}

\newcommand{\absetzungskopf}[1]{\par\vspace{6mm}\noindent #1\par\vspace{1.5mm}}

\newcommand{\absetzungstabelle}[1]{\par\noindent%
  \begin{tabularx}{\linewidth}{|>{\raggedright\arraybackslash}X|B|}\hline#1\\\hline\end{tabularx}\par}

\newcommand{\positionsdetail}[1]{\newline{\normalfont\small #1}}
\newcommand{\absetzungsdetail}[1]{\newline{\small #1}}

\newcommand{\positionszeile}[6]{#1 & \textbf{#2} & #3 & #4 & #5 & #6}

\newcommand{\summenzeile}[2]{#1 & #2}

\newcommand{\abschlusszeile}[2]{\\\hline\textbf{#1} & \textbf{#2}}

\newcommand{\bankblock}{}

\newfontfamily{\logoschrift}{texgyreheros}[
  Extension   = .otf,
  UprightFont = *-regular,
  BoldFont    = *-bold,
]
\definecolor{logodunkel}{HTML}{1B2A4A}
\definecolor{logohell}{HTML}{4FB3F0}
\definecolor{logomittel}{HTML}{2F7FD8}
\definecolor{logofalte}{HTML}{1F5FB0}
\newlength{\logohoehe}\setlength{\logohoehe}{12mm}
\newlength{\versalhoehe}
\newcommand{\logo}{\resizebox{!}{\logohoehe}{\begin{tikzpicture}
  \shade[top color=logohell, bottom color=logomittel, rounded corners=3pt]
    (0,0) -- (0,2.6) -- (0.8,2.6) -- (0.8,0) -- cycle;
  \shade[top color=logohell, bottom color=logomittel]
    (0.8,2.6) -- (1.6,1.5) -- (1.6,0.6) -- (0.8,1.7) -- cycle;
  \shade[top color=logohell, bottom color=logomittel]
    (1.6,1.5) -- (2.4,2.6) -- (2.4,1.7) -- (1.6,0.6) -- cycle;
  \shade[top color=logohell, bottom color=logomittel, rounded corners=3pt]
    (2.4,0) -- (2.4,2.6) -- (3.2,2.6) -- (3.2,0) -- cycle;
  \fill[logofalte] (0.8,2.6) -- (0.8,1.7) -- (1.15,2.12) -- cycle;
  \fill[logofalte] (2.4,2.6) -- (2.4,1.7) -- (2.05,2.12) -- cycle;
  \node[anchor=south west, inner sep=0, font=\logoschrift\fontsize{40}{40}\selectfont\bfseries]
    at (4.0,0.95) {\textcolor{logodunkel}{MUSTER}\textcolor{logohell}{IT}};
  \node[anchor=south west, inner sep=0, text=logodunkel,
    font=\logoschrift\fontsize{13}{13}\selectfont\addfontfeature{LetterSpace=30}]
    at (4.05,0.05) {CONSULTING};
\end{tikzpicture}}}

\begin{document}

\settoheight{\versalhoehe}{\LARGE\bfseries M}%
\noindent\raisebox{\dimexpr\versalhoehe-\height\relax}[0pt][0pt]{\logo}%
\hfill
\begin{tabular}[t]{@{}r@{\hspace{3mm}}l@{}}
  \multicolumn{2}{@{}r@{}}{\LARGE\bfseries Rechnung}\\[2mm]
  \kopfangabe{Rechnungsnr.:}{R26-042}
  \kopfangabe{Kundennr.:}{}
  \kopfangabe{Unsere Auftragsnr.:}{}
  \kopfangabe{Datum:}{14.~August 2026}
  \kopfangabe{Leistungszeitraum:}{Juli 2026}
  \kopfangabe{Vereinnahmt am:}{}
  \kopfangabe{Ihre Referenz:}{BST-2026-0815}
  \kopfangabe{Zahlbar bis:}{28.~August 2026}
\end{tabular}

\vspace{12mm}

\begin{minipage}[t]{0.5\linewidth}\raggedright
  Beispiel Software GmbH\zeilenwechsel{}Industriestraße 1\zeilenwechsel{}90402 Nürnberg\ifblank{}{}{\\[2mm]USt-IdNr.: }
\end{minipage}%
\hfill
\begin{minipage}[t]{0.45\linewidth}\raggedleft
  Muster IT-Consulting UG (haftungsbeschränkt)\ifblank{Dr. Erika Musterfrau}{\zeilenwechsel{}Dr. Erika Musterfrau}{}\zeilenwechsel{}Beispielweg 42\zeilenwechsel{}97070 Würzburg\ifblank{+49 931 1234567}{}{\zeilenwechsel{}Tel.: +49 931 1234567}\zeilenwechsel{}rechnung@muster-it-consulting.example
\end{minipage}

\vspace{12mm}

\noindent
\begin{tabularx}{\linewidth}{|>{\centering\arraybackslash}p{7mm}|>{\raggedright\arraybackslash}X|>{\raggedleft\arraybackslash}p{13mm}|>{\raggedright\arraybackslash}p{15mm}|>{\raggedleft\arraybackslash}p{22mm}|B|}
  \hline
  \textbf{Pos.} & \textbf{Bezeichnung} & \textbf{Menge} & \textbf{Einheit} & \textbf{Einzelpreis} & \textbf{Gesamt}\\
  \hline
\positionszeile{1}{IT-Beratungsleistungen\positionsdetail{Rahmenvertrag vom 01.02.2026}}{86,5}{Std.}{95,00~€}{8.217,50~€} \positionstrenner{}
\positionszeile{2}{Fahrtkostenpauschale\positionsdetail{Juli 2026}}{1}{pauschal}{120,00~€}{120,00~€}\\
  \hline
\end{tabularx}
\par\nointerlineskip\noindent
\begin{tabularx}{\linewidth}{|>{\raggedright\arraybackslash}X|B|}
  \summenzeile{Nettobetrag}{8.337,50~€} \summentrenner{}
  \summenzeile{zzgl. 19~\% USt.}{1.584,13~€}
  \abschlusszeile{Gesamtbetrag}{9.921,63~€}\\
  \hline
\end{tabularx}







\bankblock{}

\vspace{8mm}

Ich bedanke mich für das Vertrauen und die angenehme Zusammenarbeit.

\vfill

{\footnotesize
\noindent
\begin{minipage}[t]{0.32\linewidth}\raggedright
  Muster IT-Consulting UG (haftungsbeschränkt)\zeilenwechsel{}Beispielweg 42\zeilenwechsel{}97070 Würzburg\ifblank{https://muster-it-consulting.example}{}{\zeilenwechsel{}https://muster-it-consulting.example}
\end{minipage}%
\hfill
\begin{minipage}[t]{0.34\linewidth}\raggedright
  USt-IdNr.: DE 123 456 788\ifblank{257/123/45678}{}{\\Steuernummer: 257/123/45678}\ifblank{Würzburg}{}{\\Registergericht: Würzburg\ifblank{HRB 12345}{}{, HRB 12345}}\ifblank{Dr. Erika Musterfrau}{}{\\Geschäftsführung: Dr. Erika Musterfrau}
\end{minipage}%
\hfill
\begin{minipage}[t]{0.28\linewidth}\raggedright
  Musterbank Würzburg\\IBAN: DE89 3704 0044 0532 0130 00\\BIC: MUSTDEFFXXX
\end{minipage}\par}

\end{document}
